% 관련 연구 및 문제 설정 문단. tectonic으로 컴파일: tectonic related_work.tex
% 2026-11-25 피드백(feedback/질문방-2026-11-25 (3)–(7).md) 검증 후 반영.
\documentclass[11pt,a4paper]{article}
\usepackage{fontspec}
\usepackage{xeCJK}
\setCJKmainfont[BoldFont={Apple SD Gothic Neo Bold},ItalicFont={AppleMyungjo},ItalicFeatures={FakeSlant=0.18}]{AppleMyungjo}
\setCJKsansfont{Apple SD Gothic Neo}
\setCJKmonofont{Apple SD Gothic Neo}
\xeCJKsetup{CJKspace=true,CJKecglue={},xCJKecglue=false}
\usepackage{amsmath,amssymb,mathtools,bm}
\usepackage[margin=25mm]{geometry}
\usepackage[hidelinks]{hyperref}
\linespread{1.35}
\setlength{\parskip}{0.5em}
\setlength{\parindent}{1em}
\newcommand{\E}{\mathbb{E}}
\renewcommand{\P}{\mathbb{P}}
\newcommand{\I}{I}
\newcommand{\Yt}{\tilde{Y}}
\newcommand{\Yh}{\hat{Y}}
\newcommand{\softmax}{\operatorname{softmax}}
\newcommand{\Cat}{\operatorname{Cat}}
\newcommand{\KL}{D_{\mathrm{KL}}}
\renewcommand{\refname}{참고문헌}

\title{관련 연구 및 문제 설정}
\author{}
\date{}

\begin{document}
\maketitle
\vspace{-3em}

\section{유한 부호 길이 영역과 결합 원천-채널 부호화}

6G의 초고신뢰 저지연 통신(URLLC)과 시각-언어 모델(VLM) 기반 엣지 추론은 짧은 블록과 엄격한 지연 제약 아래에서 동작한다. 그래서 부호 길이가 무한하다고 가정하는 Shannon의 점근적 결과만으로는 성능 한계를 정확히 규명할 수 없다. 원천 길이 $k$와 부호 길이 $n$이 유한한 영역에서는 비점근 정보이론으로 달성 가능성 한계(achievability bound)와 역 한계(converse bound)를 구한다 \cite{gao2025fbl}. 대표적으로 Kostina와 Verdú는 $\mathsf{d}$-기울인 정보(d-tilted information) $\jmath_S(s,d)$를 정의했다 \cite{kostina2012, kostina2013}. 그 평균인 율-왜곡 함수 $R(d)=\E[\jmath_S(S,d)]$와 분산인 원천 분산도(source dispersion) $\mathcal{V}(d)=\operatorname{Var}[\jmath_S(S,d)]$로 유한 길이 손실 압축과 손실 결합 원천-채널 부호화(JSCC)의 한계를 2차 항까지 유도했다. $R(d)$와 최적 테스트 채널은 Blahut--Arimoto(BA) 알고리즘으로 계산할 수 있다 \cite{blahut1972, arimoto1972}.

점근 영역에서는 원천-채널 분리 정리에 따라 분리 부호화(SSCC)로도 최적에 도달한다. 그러나 유한 길이에서는 분리 설계가 일반적으로 2차 항 손실을 입는다. 정규화한 채널 정보 밀도를 $I$, 정규화한 $\mathsf{d}$-기울인 정보를 $J$라 하자. 결합 부호화는 사건 $\{I>J\}$만 일어나면 원천을 왜곡 $d$ 이내로 복원한다. 반면 분리 부호화는 두 부호 사이의 중간 율 $r$에 대해 사건 $\{I>r\}\cap\{J<r\}$이 일어나야 한다. 뒤의 사건은 앞 사건의 진부분집합이므로, 같은 오류 확률에서 달성 가능한 부호율이 낮아진다. 다만 원천 분산도와 채널 분산도 가운데 하나라도 0이면 분리 부호화로도 결합 설계와 같은 분산도를 얻는다 \cite{gao2025fbl}.

딥러닝 기반 JSCC는 신경망 부호기와 복호기를 종단 간(end-to-end)으로 학습한다. 이렇게 이미지, 비디오, 의미 특징을 채널 심볼로 직접 사상한다. 초기 형태는 대역폭 압축비, 곧 채널 심볼 수와 원천 차원의 비, 그리고 학습 SNR을 고정하고 복원 왜곡(MSE)을 최소화했다 \cite{bourtsoulatze2019}. 이 방식은 채널 상태가 나빠질 때 분리 설계의 절벽 효과(cliff effect) 없이 성능이 점진적으로 저하되고, 학습 때와 다른 SNR에서도 강인함을 실험으로 보였다 \cite{bourtsoulatze2019}. 또한 채널 출력 되먹임을 도입한 JSCC는 전송 지연을 줄인다고 보고되었다 \cite{gao2025fbl}. 이후 채널 조건에 따라 압축률과 등화 여부를 적응적으로 정하는 방식으로 발전했다. Seon 등의 방식에서는 송신 측 DRL 에이전트가 추정 SNR을 바탕으로 압축률을 고르고, 수신 측 DRL 에이전트가 VAE 등화 수행 여부를 고른다 \cite{seon2025}. 두 에이전트는 추정 SNR과 짧은 참조 부호로 이루어진 부분 채널 상태 정보(CSI)를 입력으로 받는다. VAE 등화기는 수신 신호의 실수부와 허수부로부터 채널 $\hat B$를 블라인드 추정하고, $\hat B^{-1}y$로 등화한다. 학습 손실에는 실제 채널이 들어가지만, 추론할 때는 MMSE와 달리 채널 행렬 없이 등화할 수 있다.

\section{과업 지향 부호화의 목적 함수}

과업 지향 부호화의 목적 함수는 정보 병목(IB, information bottleneck)으로 정식화할 수 있다 \cite{tishby1999}. 부호기는 입력 $X$만 보고 표현 $T$를 만들므로 마르코프 사슬 $Y - X - T$가 성립한다. 여기서 $Y$는 과업 변수다.
\begin{equation}
\min_{p(t\mid x)}\; \I(X;T) - \beta\, \I(T;Y)
\label{eq:ib}
\end{equation}
$\I(X;T)$는 $T$가 $X$에 대해 담는 정보량, 곧 전송해야 할 율이다. $\I(T;Y)$는 $T$가 $Y$에 대해 담는 관련 정보, 곧 과업 성능이다. 두 항은 서로 상충하며, 라그랑주 승수 $\beta$가 그 균형을 정한다. $\I(T;Y)\le \I(X;T)$이므로 $\beta\le1$이면 $T$가 $X$와 독립인 자명한 해가 최적이고, 실제로는 어떤 임계값 $\beta_c\ge1$ 아래에서 율이 0이 된다. $\beta$가 커질수록 $T$는 데이터 처리 부등식의 상한 $\I(T;Y)\le \I(X;Y)$에 도달하는 최소 충분 통계량에 다가간다.

식~\eqref{eq:ib}는 다음 간접 율-왜곡 문제와 직접 연결된다.
\begin{equation}
\min_{p(t\mid x)}\; \I(X;T) \quad \text{s.t.} \quad \E\bigl[d(Y,\Yh)\bigr] \le D
\label{eq:ird}
\end{equation}
부호기가 보지 못한 $Y$를 복원하므로 간접(indirect) 문제라 부르며, 최적화 대상은 부호기, 곧 테스트 채널 $p(t\mid x)$다. 분류 과업에서 $Y$는 범주형 변수이므로, 왜곡 척도의 대표적인 선택은 0-1(해밍) 왜곡과 로그 손실이다. 0-1 왜곡 $d(y,\hat y)=\mathbf{1}\{y\ne\hat y\}$의 기댓값은 오류 확률 $\P(\Yh\ne Y)$다 \cite{enttsel2026}. 로그 손실에서는 복원값이 라벨이 아니라 $Y$ 위의 확률분포 $\Yh=q(\cdot\mid T)$이며, 그 기댓값은 다음과 같이 분해된다.
\begin{equation}
\E\bigl[-\log q(Y\mid T)\bigr] = H(Y\mid T) + \E\bigl[\KL\bigl(p(\cdot\mid T)\,\big\|\,q(\cdot\mid T)\bigr)\bigr]
\end{equation}
따라서 최솟값은 $q=p(\cdot\mid t)$일 때의 $H(Y\mid T)=H(Y)-\I(T;Y)$다. 곧 로그 손실 아래에서 식~\eqref{eq:ird}의 왜곡 제약은 $\I(T;Y)\ge H(Y)-D$와 같고, 이 제약 문제의 라그랑주 형태가 식~\eqref{eq:ib}이다. $q$를 신경망 분류기로 두면 이 기댓값은 분류 모델 학습에 쓰는 교차 엔트로피 손실이다.

\section{softmax 출력과 압축 전략}

softmax 기반 분류 모델은 로짓 $f(x)\in\mathbb{R}^{K}$로부터 사후 확률 벡터를 출력한다. 여기서 $K$는 클래스 수다.
\begin{equation}
V = \softmax\bigl[f(X)\bigr]
\label{eq:v}
\end{equation}
softmax는 압축 연산이 아니다. softmax가 잃는 것은 모든 로짓에 같은 상수를 더한 차이뿐이고, 이 상수는 $Y$와 무관하다. 모델의 사후 확률을 참으로 간주하면, 과업 모델을 고정했을 때 $V$는 $Y$에 대한 $X$의 베이즈 충분 통계량이다 \cite{enttsel2026}. 이와 별개로, 확정 라벨 $\Yt=\arg\max V$는 $V$의 함수이므로 $Y - X - V - \Yt$가 마르코프 사슬을 이루고, 데이터 처리 부등식에 따라 $\I(Y;V)\ge \I(Y;\Yt)$다. 이 부등식은 충분성 가정 없이 성립한다.

Enttsel과 Corlay는 과업 모델에 넣기 전에 $X$를 압축하는 압축 후 추정(C\&E)과 함께, 모델 출력 뒤에서 압축하는 전략들을 분석했다 \cite{enttsel2026}. 이 글에서는 뒤의 전략들을 세 갈래로 묶어 정리한다.

첫째, 추정 후 압축(E\&C)은 확정 라벨 $\Yt$를 낸 뒤 압축한다. 무손실 전송에는 $H(\Yt)$ 비트면 충분하며, 이 값은 $\log_2 K$ 이하다. 그보다 낮은 율에서의 한계는 BA 알고리즘으로 계산한다. 이를 개선한 간접 추정 후 압축(iE\&C)은 압축기가 $\Yt$ 자체가 아니라 진짜 라벨 $Y$를 복원하도록 다음 유효 왜곡을 사용한다. 둘째 등호는 $d$가 0-1 왜곡일 때 성립한다.
\begin{equation}
\hat d(\tilde y,\hat y) = \E\bigl[d(Y,\hat y)\mid \Yt=\tilde y\bigr] = 1 - \P\bigl(Y=\hat y\mid \Yt=\tilde y\bigr)
\end{equation}
그 결과 iE\&C는 E\&C와 같거나 더 낮은 율 한계를 준다. 다만 두 방법 모두 입력이 $\Yt$뿐이어서 $V$에 담긴 모델의 확신도 정보를 버린다.

둘째, 모델 인지 율-왜곡(TMRD)은 $V$ 전체를 원천으로 삼아 간접 율-왜곡 문제를 푼다.
\begin{equation}
\min_{p(\hat y\mid v)}\; \I(V;\Yh) \quad \text{s.t.} \quad \E\bigl[d(Y,\Yh)\bigr]\le D
\end{equation}
이 한계는 고정된 과업 모델로 얻을 수 있는 최적의 율-왜곡 성능, 곧 최적 한계다. 그러나 $V$가 $K$차원 연속 벡터여서 BA 알고리즘으로 풀 수 없으므로, TMRD는 계산할 수 없는 기준선으로 남는다.

셋째, 그 계산 가능한 대안으로 제시된 것이 표본 추출 후 전송(S\&C)이다. 확정 라벨 대신 사후 분포에서 표본을 하나 뽑는다.
\begin{equation}
(\Yh\mid X=x) \sim \Cat\bigl(\softmax[\lambda f(x)]\bigr)
\end{equation}
원문은 이를 채널 시뮬레이션(상대 엔트로피 부호화)의 한 예로 보고, 율을 $H(\Yh)$가 아니라 $\I(X;\Yh)$로 잡는다. 역온도 $\lambda\to\infty$이면 무손실 동작점 $(H(\Yt), D_{\mathrm{TM}})$을 되찾고, $\lambda\to0$이면 $\Yh$가 $X$와 독립이 되어 율이 0이 된다. 여기서 $D_{\mathrm{TM}}$은 과업 모델 자체의 오류율이다. 이 방식은 확신도를 표본의 무작위성으로 반영하지만 명시적인 최적성 보장은 없다.

\section{분할 지점}

위 전략들은 모두 과업 모델 전체가 송신기에서 실행된다고 가정한다. 송신기의 메모리와 연산이 제한되면 모델을 분할 지점에서 나눠야 한다. 앞쪽의 과업 전단(front end)은 송신기에서 실행해 중간 특징을 만들고, 이 특징을 압축하거나 JSCC로 전송한다. 뒤쪽의 과업 후단(back end)과 softmax는 수신기에서 실행한다.

Harell 등은 최적 부호화를 가정하고, 특징의 근사값을 과업 후단에 넣은 출력이 과업 모델의 상(image) 안에 있다는 조건 아래에서, 분할 지점이 최소 달성 율에 영향을 주지 않음을 증명했다 \cite{harell2025}. 분류 문제에서는 이 조건이 자명하게 성립한다. 그러나 이 결과는 이론적 최적의 경우에 관한 것이며, 원문도 실제 코덱은 최적이 아니어서 입력을 직접 부호화하는 방식과 모델 분할 방식 사이에 큰 차이가 날 수 있다고 지적한다. 실제 시스템에서 분할 지점은 다음 두 관점에서 성능을 좌우한다.

첫째, 정보이론 관점에서 Enttsel과 Corlay는 배치된 과업 모델을 명시적으로 반영하면 압축 대상에 따라 한계가 달라짐을 보였다 \cite{enttsel2026}.
\begin{equation}
R_Y(D) \le R_X(D) \le R_{\Yt}^{\mathrm{iE\&C}}(D) \le R_{\Yt}^{\mathrm{E\&C}}(D)
\label{eq:order}
\end{equation}
여기서 $R_Y$는 $Y$를 직접 관측한다고 가정한 하한, $R_X$는 $X$로부터의 간접 율-왜곡 함수, 뒤의 두 항은 확정 라벨 $\Yt$를 압축하는 iE\&C와 E\&C의 한계다. 원문은 TMRD 한계가 하한 $R_Y$와 과업 모델 의존 한계들 사이에 놓일 것으로 보지만, 이는 정성적 주장이다. 같은 ResNet-50을 서로 다른 두 지점에서 분할한 기존 방식들도 서로 다른 율-정확도 곡선을 보였다. 송신기 모델 크기는 Layer 4의 Block 0에서 나누면 28\,MB, Layer 3의 Block 4에서 나누면 12.16\,MB다. 저자들은 이론 한계와의 율 격차가 주로 송신기 복잡도 제약과 그에 따른 모델 분할 때문이라고 결론지었다.

둘째, 시스템 관점에서 분할 지점은 전송량과 지연을 직접 결정한다. ResNet 계열에서는 거의 모든 중간 특징이 입력보다 큰 데이터 증폭 효과가 나타나, 분할 지점만 바꿔서는 추론 지연을 줄이기 어렵다 \cite{shao2020tradeoff}. Edgent에서는 분할 자체가 정확도를 바꾸지 않는다. 대신 지연 제약 아래에서 분할 지점이 지연을 정하고, 그에 따라 고를 수 있는 조기 종료 지점과 정확도가 정해진다 \cite{li2020edgent}. 페이딩 채널에서 Yan 등은 분할 지점 선택을 최적 정지 문제로 정식화했다. 그리고 단말에 올린 층 수가 주어졌을 때, 관측 SNR이 단계별 임계값 이상이 되는 첫 지점에서 나누는 임계값 정책이 기대 에너지-시간 비용을 최소화함을 증명했다 \cite{yan2021split}. VLA 모델에서도 대역폭이 10\,MB/s에서 1\,MB/s로 떨어지면 기존 최적 분할 지점의 전송 지연이 9.9\,ms에서 99.6\,ms로 늘어난다. 다른 지점으로 옮기면 24.9\,ms로 줄어, 최적 분할 지점이 이동한다는 측정 결과가 보고되었다 \cite{zheng2026roboecc}.

따라서 송신기 복잡도, 고정된 과업 모델, 시변 채널이라는 현실적 제약 아래에서 분할 지점은 율, 정확도, 지연을 함께 결정하는 핵심 설계 변수다. 이 연구는 이 제약 아래에서 부호기의 분할 지점과 구조를 정하는 문제를 다룬다.

\begin{thebibliography}{99}\small
\bibitem{gao2025fbl} J. Gao, S. Chen, Y. Wu, L. Liu, G. Caire, H. V. Poor, and W. Zhang, ``Finite-blocklength information theory,'' arXiv:2504.07743, 2025.
\bibitem{kostina2012} V. Kostina and S. Verdú, ``Fixed-length lossy compression in the finite blocklength regime,'' \emph{IEEE Trans. Inf. Theory}, vol.~58, no.~6, pp.~3309--3338, 2012.
\bibitem{kostina2013} V. Kostina and S. Verdú, ``Lossy joint source-channel coding in the finite blocklength regime,'' \emph{IEEE Trans. Inf. Theory}, vol.~59, no.~5, pp.~2545--2575, 2013.
\bibitem{blahut1972} R. E. Blahut, ``Computation of channel capacity and rate-distortion functions,'' \emph{IEEE Trans. Inf. Theory}, vol.~18, no.~4, pp.~460--473, 1972.
\bibitem{arimoto1972} S. Arimoto, ``An algorithm for computing the capacity of arbitrary discrete memoryless channels,'' \emph{IEEE Trans. Inf. Theory}, vol.~18, no.~1, pp.~14--20, 1972.
\bibitem{bourtsoulatze2019} E. Bourtsoulatze, D. B. Kurka, and D. Gündüz, ``Deep joint source-channel coding for wireless image transmission,'' \emph{IEEE Trans. Cogn. Commun. Netw.}, vol.~5, no.~3, pp.~567--579, 2019.
\bibitem{seon2025} J. Seon, S. Lee, S. H. Kim, Y. G. Sun, H. Seo, D. I. Kim, and J. Y. Kim, ``Deep reinforced feature compression and channel equalization for semantic communications,'' \emph{IEEE Trans. Cogn. Commun. Netw.}, vol.~11, no.~6, pp.~3655--3668, 2025.
\bibitem{tishby1999} N. Tishby, F. C. Pereira, and W. Bialek, ``The information bottleneck method,'' in \emph{Proc. Annu. Allerton Conf. Commun., Control, Comput.}, 1999, pp.~368--377.
\bibitem{enttsel2026} A. Enttsel and V. Corlay, ``Model-aware rate-distortion limits for task-oriented source coding,'' arXiv:2602.12866, 2026.
\bibitem{harell2025} A. Harell, Y. Foroutan, N. Ahuja, P. Datta, B. Kanzariya, V. S. Somayazulu, O. Tickoo, A. de Andrade, and I. V. Bajić, ``Rate-distortion theory in coding for machines and its applications,'' \emph{IEEE Trans. Pattern Anal. Mach. Intell.}, vol.~47, no.~7, pp.~5501--5519, 2025.
\bibitem{shao2020tradeoff} J. Shao and J. Zhang, ``Communication-computation trade-off in resource-constrained edge inference,'' \emph{IEEE Commun. Mag.}, vol.~58, no.~12, pp.~20--26, 2020.
\bibitem{li2020edgent} E. Li, L. Zeng, Z. Zhou, and X. Chen, ``Edge AI: On-demand accelerating deep neural network inference via edge computing,'' \emph{IEEE Trans. Wireless Commun.}, vol.~19, no.~1, pp.~447--457, 2020.
\bibitem{yan2021split} J. Yan, S. Bi, and Y.-J. A. Zhang, ``Optimal model placement and online model splitting for device-edge co-inference,'' arXiv:2105.13618, 2021.
\bibitem{zheng2026roboecc} Z. Zheng \emph{et al.}, ``RoboECC: Multi-factor-aware edge-cloud collaborative deployment for VLA models,'' arXiv:2603.20711, 2026.
\end{thebibliography}

\end{document}
